\chapter{\texttt{core.BSM} --- Beyond-the-Standard-Model extensions}
\label{ch:core-bsm}

This chapter covers the two supported BSM extensions: non-standard
interactions (NSI) in the 3-flavour sector, and a fourth light sterile
neutrino (the ``3+1'' scheme). Both reuse as much as possible of the
infrastructure from \cref{ch:core-common}: NSI only modifies $H_{\rm mat}$,
reusing whichever \pyclass{PMNS} object is otherwise in use, and the
sterile sector extends the factorisation $U=O\,U_{\rm red}$ to
$4\times4$. There is no separate \texttt{core.BSM} Hamiltonian-assembly
module: \pyfunc{hamiltonian\_reduced} and \pyfunc{hamiltonian\_flavour}
(\cref{ch:core-common}, \texttt{hamiltonian.py} section) already handle NSI,
the sterile extension, or both simultaneously.

\section{\texttt{bsm\_nsi.py} --- non-standard interactions}

\modulehead{core/BSM/bsm\_nsi.py}{parametric configuration of NSI couplings
and construction of the $\epsilon$ matrix.}

\begin{theorybox}
NSI are effective four-fermion operators that modify coherent forward
scattering beyond the standard MSW potential \parencite{Grossman1995}. The
matter Hamiltonian in the flavour basis generalises to
\begin{equation}
  \keyresult{H_{\rm mat}^{\rm NSI} = V\left(\begin{pmatrix}1&0&0\\0&0&0\\0&0&0\end{pmatrix} + \epsilon\right)},
  \qquad V = \pm\sqrt2\,G_F\,n_e\,L_{\rm scale},
\end{equation}
with $\epsilon$ a dimensionless $3\times3$ Hermitian matrix,
\begin{equation}
  \epsilon = \begin{pmatrix}
    \epsilon_{ee} & \epsilon_{e\mu} & \epsilon_{e\tau} \\
    \epsilon_{e\mu}^* & \epsilon_{\mu\mu} & \epsilon_{\mu\tau} \\
    \epsilon_{e\tau}^* & \epsilon_{\mu\tau}^* & \epsilon_{\tau\tau}
  \end{pmatrix}.
\end{equation}
Diagonal entries are real; off-diagonal entries are generally complex. The
Standard Model limit is $\epsilon=0$. Only the traceless part of $\epsilon$
matters for oscillations (an overall trace is an irrelevant global phase),
so fixing $\epsilon_{\mu\mu}=0$ is a valid convention for the diagonal
sector without loss of physical generality.

\textbf{No NSI-specific mixing matrix.} \pyclass{NSIConfig} touches only
$H_{\rm mat}$: it does not define, or require, any dedicated \pyclass{PMNS}
subclass. If the sterile extension is not also active, NSI propagation
reuses the plain 3-flavour \pyclass{PMNS\_SM} object from \cref{ch:core-sm}
completely unchanged -- the same $U$, $O$, $U_{\rm red}$ as the pure
Standard Model case. \pyclass{OscillationParameters.nsi} is simply an
optional, orthogonal attribute alongside \code{pmns}
(\cref{ch:core-common}, \texttt{oscillation.py} section): which mixing
scenario is active and whether NSI is active are two independent choices,
and every combination (SM+NSI, sterile alone, sterile+NSI) is valid.
\end{theorybox}

\begin{theorybox}
\subsection*{Transforming the matter Hamiltonian: Standard Model versus NSI}
In the plain Standard Model, the reduced-basis matter term needs no
rotation at all: $H_{\rm mat}^{\rm red}=H_{\rm mat}^{\rm flavour}=\diag(V,0,\dots,0)$,
because $O\,H_{\rm mat}\,O^\dagger=H_{\rm mat}$ exactly (\cref{ch:core-sm}) --
the matter term is identical in the flavour and reduced bases, so it is
simply added to the reduced kinetic term as-is.

When NSI are active this shortcut disappears: as shown below,
$O\,H_{\rm mat}^{\rm NSI}\,O^\dagger\neq H_{\rm mat}^{\rm NSI}$ in general.
\pyfunc{hamiltonian\_matter\_reduced} (\cref{ch:core-common}) therefore
takes the opposite route: it builds $H_{\rm mat}^{\rm NSI}$ directly in the
\emph{flavour} basis as above, and then rotates it explicitly \emph{down}
into the reduced basis,
\begin{equation}
  \keyresult{H_{\rm mat}^{\rm NSI,\,red} = O^\dagger\,H_{\rm mat}^{\rm NSI,\,flavour}\,O},
\end{equation}
the exact inverse of the flavour-basis dressing $O(\cdot)O^\dagger$
(\cref{ch:core-common}). Because the transformation is performed
explicitly here -- rather than assumed to be a no-op, as in the SM case
above -- the identity
$H_{\rm flavour}=\text{\pyfunc{pmns.flavour\_basis}}(H_{\rm red})$ holds
exactly for \emph{arbitrary} $\epsilon$, not only for the
diagonal-uniform Standard Model limit.
\end{theorybox}

\begin{notebox}
\textbf{Adjoint, not transpose.} $O$ is generally complex ($\Delta$ carries
$e^{i\delta}$), so $O^\dagger=(O^{*})^{T}$ and the plain transpose $O^{T}$
differ whenever $\delta\neq0$. The rotation above \textbf{must} use the
conjugate transpose (\code{O.conj().transpose(-2,-1)} in
\pyfunc{hamiltonian\_matter\_reduced}) to preserve Hermiticity: for
Hermitian $\epsilon$ (and hence Hermitian $H_{\rm mat}^{\rm NSI}$),
$O^\dagger\,H\,O$ is again Hermitian for any unitary $O$, whereas $O^{T}\,H\,O$
is not, in general, even Hermitian -- let alone the physically correct
transformed operator. This mirrors the identical requirement already
established for the kinetic term (\cref{ch:core-common},
\texttt{pmns.py}/\texttt{hamiltonian.py} sections): the Hermitian conjugate
is the only convention consistent with a unitary change of basis for an
operator; the plain transpose coincides with it only by accident, when the
matrix involved happens to be real.
\end{notebox}

\begin{notebox}
Since $H_{\rm mat}^{\rm NSI}$ is no longer proportional to $\diag(V,0,0)$,
the commutation property $O\,H_{\rm mat}\,O^\dagger=H_{\rm mat}$ assumed by
the general factorisation \cref{eq:H-factorization} \textbf{no longer holds
in general} once $\epsilon$ has entries outside the $ee$ block. The code
does not rely on it being true -- as shown above,
\pyfunc{hamiltonian\_matter\_reduced} builds $H_{\rm mat}^{\rm NSI}$ in the
flavour basis and rotates it explicitly, so the result is exact for any
$\epsilon$. What \emph{is} lost is the Standard Model \emph{simplification}
(\cref{ch:core-sm}) that $H_{\rm red}$ needs no reference to
$\theta_{23}/\delta$ and can be built with a real $U_{\rm red}$: for
generic NSI, $H_{\rm red}$ is genuinely complex and depends on every
mixing parameter, so the closed-form spectral formulas used elsewhere in
the project no longer apply, and the relevant diagonalisation (e.g.\ in
\texttt{medium.solar.probability}) is performed numerically instead, via
\code{torch.linalg.eigh}.
\end{notebox}

\begin{implbox}
\begin{itemize}
  \item \pyclass{NSIConfig} (frozen dataclass): stores the six independent
    parameters of $\epsilon$ (3 real diagonal + 3 complex off-diagonal) and
    exposes \pyfunc{epsilon\_tensor\_base()} (the raw complex $3\times3$
    tensor), \pyfunc{epsilon\_tensor(n\_flavours=...)} (embedded for any
    flavour count, zero-padded for the sterile rows/columns), and
    \pyfunc{is\_sm\_limit}.
  \item \pyfunc{NSIConfig.from\_preset(name)} builds named configurations
    from \code{tpeanuts.config.presets.NSI\_PRESETS}
    (\cref{sec:nsi-presets}), each with its physical justification
    documented.
  \item The $\epsilon$ tensor is read via
    \pyfunc{oscillation.nsi.epsilon\_tensor(...)}, passed through
    \pyfunc{pmns.select\_antinu} ($\epsilon\to\epsilon^{*}$ for
    antineutrinos, mirroring the $U\to U^{*}$ convention), and injected
    directly into \pyfunc{hamiltonian\_matter\_reduced}
    (\cref{ch:core-common}) -- the single Hamiltonian-assembly function
    handling this case; there is no separate \texttt{core.BSM} builder.
\end{itemize}
\end{implbox}

\begin{refbox}
Original NSI proposal: \cite{Grossman1995}. Model-independent bounds:
\cite{BiggioBlennowFM2009}. LMA-Dark degeneracy in the global fit:
\cite{Esteban2018LMADark}. IceCube DeepCore bounds: \cite{IceCube2022NSI}.
\end{refbox}

\section{NSI presets}
\label{sec:nsi-presets}

\modulehead{tpeanuts/config/presets.py}{\code{NSI\_PRESETS}: named,
literature-justified $\epsilon$ configurations, built via
\pyfunc{NSIConfig.from\_preset}.}

\begin{implbox}
Every preset below sets a handful of $\epsilon$ parameters to a
representative, literature-motivated value (all other components default
to zero) rather than to a unique global best fit: current data leaves large
NSI degeneracies, so no single preferred point exists. Their purpose is
threefold: (1) give the validation notebooks and tests physically motivated
$\epsilon$ values instead of arbitrary numbers, covering diagonal-only,
off-diagonal, single- and multi-parameter regimes; (2) reproduce, and stay
traceable to, specific published sensitivity/bound studies; and
(3) let users reach a standard benchmark scenario (e.g.\ the LMA-Dark
degeneracy) by name instead of re-deriving its $\epsilon$ values.

\begin{center}
\small
\begin{tabular}{@{}ll@{}}
\toprule
Preset & $\epsilon$ values and origin \\
\midrule
\parbox[t]{3.2cm}{\code{sm\_\allowbreak{}no\_\allowbreak{}nsi}} &
  \parbox[t]{9.4cm}{All $\epsilon=0$. Standard Model limit; equivalent to
  passing \code{nsi=None}.} \\[10pt]
\parbox[t]{3.2cm}{\code{nsi\_\allowbreak{}diagonal\_\allowbreak{}biggio2009}} &
  \parbox[t]{9.4cm}{$\epsilon_{ee}=0.30$, $\epsilon_{\tau\tau}=0.15$.
  Diagonal-only benchmark within the model-independent 90\% CL bounds of
  \textcite{BiggioBlennowFM2009} ($|\epsilon_{ee}|<0.50$,
  $|\epsilon_{\tau\tau}|<0.33$).} \\[18pt]
\parbox[t]{3.2cm}{\code{nsi\_\allowbreak{}lma\_\allowbreak{}dark\_\allowbreak{}esteban2018}} &
  \parbox[t]{9.4cm}{$\epsilon_{ee}=-2.0$. The LMA-Dark degenerate solar
  solution of \textcite{Esteban2018LMADark}: this large negative
  $\epsilon_{ee}$ flips the sign of the effective MSW potential, mimicking
  $\theta_{12}\to\pi/2-\theta_{12}\approx56.6^\circ$ (paired with the
  \code{"\_LMA\_DARK\_NUFIT52\_NO"} oscillation preset,
  \cref{ch:core-common}).} \\[24pt]
\parbox[t]{3.2cm}{\code{nsi\_\allowbreak{}offdiag\_\allowbreak{}icecube2022}} &
  \parbox[t]{9.4cm}{$\epsilon_{\mu\tau}=0.005$, $\epsilon_{\tau\tau}=0.015$.
  Sits near the 90\% CL boundary of the IceCube DeepCore atmospheric
  analysis \parencite{IceCube2022NSI}
  ($|\epsilon_{\mu\tau}|<0.0060$, $|\epsilon_{\tau\tau}-\epsilon_{\mu\mu}|<0.018$).} \\[18pt]
\parbox[t]{3.2cm}{\code{nsi\_\allowbreak{}dune\_\allowbreak{}etau}} &
  \parbox[t]{9.4cm}{$\epsilon_{e\tau}=0.15$ (real). Representative DUNE
  Near Detector sensitivity benchmark \parencite{DUNE2021BSM}: this
  flavour-changing coupling produces the strongest degeneracies with
  $\delta_{CP}$, the mass ordering, and $\theta_{23}$.} \\[18pt]
\parbox[t]{3.2cm}{\code{nsi\_\allowbreak{}hyperk\_\allowbreak{}etau}} &
  \parbox[t]{9.4cm}{$\epsilon_{ee}=0.20$, $\epsilon_{e\tau}=0.10$ (real).
  Combines a diagonal MSW correction with a flavour-changing coupling,
  representative of Hyper-Kamiokande sensitivity studies
  \parencite{HyperK2015Proposal}.} \\[18pt]
\parbox[t]{3.2cm}{\code{nsi\_\allowbreak{}globalfit\_\allowbreak{}esteban2018}} &
  \parbox[t]{9.4cm}{$\epsilon_{ee}=0.25$, $\epsilon_{e\tau}=0.12$,
  $\epsilon_{\mu\tau}=0.01$, $\epsilon_{\tau\tau}=0.08$. Multi-parameter
  benchmark within the allowed region of the global oscillation analysis of
  \textcite{Esteban2018LMADark}; not a unique best fit, since significant
  correlations remain.} \\[22pt]
\parbox[t]{3.2cm}{\code{nsi\_\allowbreak{}ee\_\allowbreak{}only}} &
  \parbox[t]{9.4cm}{$\epsilon_{ee}=0.30$. Simplest possible NSI scenario:
  only the effective MSW potential is modified, with no flavour-changing
  coupling \parencite{BiggioBlennowFM2009}.} \\[10pt]
\parbox[t]{3.2cm}{\code{nsi\_\allowbreak{}mutau\_\allowbreak{}only}} &
  \parbox[t]{9.4cm}{$\epsilon_{\mu\tau}=0.01$ (real). Atmospheric
  $\mu$--$\tau$ benchmark representative of IceCube/KM3NeT sensitivity to
  flavour-changing matter interactions \parencite{Coloma2017IceCubeNSI}.} \\[14pt]
\bottomrule
\end{tabular}
\end{center}
\end{implbox}

\begin{refbox}
\cite{BiggioBlennowFM2009,Esteban2018LMADark,IceCube2022NSI,DUNE2021BSM,HyperK2015Proposal,Coloma2017IceCubeNSI}.
\end{refbox}

\section{\texttt{bsm\_sterile.py} --- the 3+1 scheme}

\modulehead{core/BSM/bsm\_sterile.py}{\pyclass{PMNS\_sterile}: $4\times4$
mixing for one additional light sterile neutrino.}

\begin{theorybox}
A fourth mass eigenstate $\nu_4$ and a fourth sterile flavour eigenstate
$\nu_s$ are added; $\nu_s$ is a gauge singlet and therefore couples to
neither $W$ nor $Z$. The charged-current potential is consequently
unaffected by the extension, so, omitting the optional neutral-current term
(the default, \cref{ch:core-common}), the matter potential in the extended
basis $(e,\mu,\tau,s)$ is
\begin{equation}
  H_{\rm mat}=\diag(V_{CC},0,0,0),
\end{equation}
the same charged-current-only term as the 3-flavour case, simply padded
with a zero sterile entry. Since $\nu_s$ also does not receive the
neutral-current potential $V_{NC}$, optionally supplying a neutron density
activates the relative term derived in \cref{ch:core-common},
\begin{equation}
  H_{\rm mat}=\diag(V_{CC},0,0,-V_{NC}),
\end{equation}
which does affect sterile-active oscillation probabilities (see the
implementation box below for how this interacts with $O_4$). The $4\times4$
mixing matrix is parametrised by adding three active--sterile rotations,
$R_{14},R_{24},R_{34}$, to the standard structure, following the
parametrisation convention of \textcite{KoppMachadoMaltoniSchwetz2013}:
\begin{equation}
  U_4 = R_{23}\,\Delta_4\,R_{24}\,R_{34}\,R_{13}\,\Delta_4^\dagger\,R_{12}\,R_{14}.
\end{equation}
Identifying $O_4\equiv R_{23}\,\Delta_4\,R_{24}\,R_{34}$
(\pyfunc{PMNS\_sterile.outer\_block}) leaves a formal reduced factor
$U_{\rm red,4}\equiv R_{13}\,\Delta_4^\dagger\,R_{12}\,R_{14}$ satisfying
$U_4=O_4\,U_{\rm red,4}$ exactly, mirroring \cref{ch:core-sm}. As in the SM
case, $\Delta_4$ cancels from the reduced \emph{kinetic term} specifically:
\begin{equation}
  \keyresult{\Delta_4^\dagger\bigl(R_{12}R_{14}\diag(k)R_{14}^\dagger R_{12}^\dagger\bigr)\Delta_4
  = R_{12}R_{14}\diag(k)R_{14}^\dagger R_{12}^\dagger},
\end{equation}
valid because neither $R_{12}$ (the $e$--$\mu$ plane) nor $R_{14}$ (the
$e$--$s$ plane) mixes with the $\tau$ index, which is the only non-trivial
index of $\Delta_4$. This is why \pyfunc{PMNS\_sterile.reduced} returns the
simplified, $\Delta_4$-free
\begin{equation}
  U_{\rm red,4}^{\rm code}=R_{13}R_{12}R_{14}
\end{equation}
rather than the formal $U_{\rm red,4}$ above: the two agree once contracted
with the kinetic term (any diagonal unitary conjugating a diagonal matrix
leaves it invariant), but $U_{\rm red,4}^{\rm code}$ is cheaper. Note that,
consequently, $O_4\,U_{\rm red,4}^{\rm code}\neq U_4$ as bare matrices --
only $U_{\rm red,4}$ (with $\Delta_4^\dagger$ included) satisfies the exact
factorisation; \pyfunc{PMNS\_sterile.reduced} returns the code-optimised
form, valid specifically for building $H_{\rm kin}$, not a literal factor
of $U_4$ (see the limitations box below). The same commutation argument as
\cref{ch:core-sm} applies to $O_4$ for the plain charged-current matter
term: $R_{23},\Delta_4,R_{24},R_{34}$ all leave the $e$ index fixed, so
$O_4\,\diag(V_{CC},0,0,0)\,O_4^\dagger=\diag(V_{CC},0,0,0)$ exactly, and the
reduced Hamiltonian is assembled and dressed back to the flavour basis
exactly as in \cref{ch:core-common,ch:core-sm}, now for $n=4$. As
established in \cref{ch:core-common} (\texttt{hamiltonian.py} section),
this commutation is specific to the charged-current-only term: $O_4$ mixes
the $\{\mu,\tau,s\}$ block and therefore does \emph{not} leave the optional
neutral-current term $\diag(0,0,0,V_{NC})$ invariant, so the reduction-trick
shortcut used above breaks whenever that term is active (see the
implementation box below).

\textbf{SM limit.} When $\theta_{14}=\theta_{24}=\theta_{34}=0$, the
$4\times4$ matrices block-diagonalise exactly into the $3\times3$ SM matrix
\emph{embedded} in the larger space, recovering every result of
\cref{ch:core-common,ch:core-sm} without approximation.

\textbf{CP-phase counting.} For $N=4$ Dirac generations there are
$(N-1)(N-2)/2=3$ physical phases: $\delta_{13}$ (the standard phase, in the
$\Delta_4$ carried through $R_{13}$), $\delta_{14}$ ($e$--$s$ sector, in
$R_{14}$), and $\delta_{24}$ ($\mu$--$s$ sector, in $R_{24}$). $R_{34}$ is
always real: the fourth possible phase can be absorbed into a field
redefinition, so sterile presets do not include \code{delta34\_deg}.
\end{theorybox}

\begin{theorybox}
\subsection*{Transforming to the reduced basis: $H_{\rm kin}$ and $H_{\rm mat}$}
Exactly as in \cref{ch:core-sm}, the reduced-basis kinetic term uses the
code's $U_{\rm red,4}^{\rm code}=R_{13}R_{12}R_{14}$ together with its
\emph{Hermitian conjugate} (never the plain transpose, for the same reason
as \cref{ch:core-common,ch:core-sm}):
\begin{equation}
  \keyresult{H_{\rm kin,4}^{\rm red} = U_{\rm red,4}^{\rm code}\,\diag(k_1,k_2,k_3,k_4)\,\bigl(U_{\rm red,4}^{\rm code}\bigr)^\dagger}.
\end{equation}
Unlike the SM's $R_{13}R_{12}$, $U_{\rm red,4}^{\rm code}$ is genuinely
complex whenever the active-sterile CP phase $\delta_{14}\neq0$ (it carries
$R_{14}$, which -- unlike $R_{12}$/$R_{13}$ -- is built with a non-trivial
phase); the Hermitian-conjugate convention is therefore not merely a safe
default here, but the only correct choice.

\textbf{Without NSI and without the neutral-current term} (the default),
the matter term needs no basis change at all:
$O_4\,\diag(V_{CC},0,0,0)\,O_4^\dagger=\diag(V_{CC},0,0,0)$ exactly (shown
above), so
\begin{equation}
  H_{\rm mat,4}^{\rm red} = H_{\rm mat,4}^{\rm flavour} = \diag(V_{CC},0,0,0),
\end{equation}
and the reduced Hamiltonian is simply
$H_{\rm red,4}=H_{\rm kin,4}^{\rm red}+H_{\rm mat,4}^{\rm red}$
(\pyfunc{hamiltonian\_reduced}, \cref{ch:core-common}), dressed back to the
flavour basis with $O_4(\cdot)O_4^\dagger$ (\pyfunc{hamiltonian\_flavour}).
This case introduces no approximation beyond \cref{ch:core-sm}: the
reduced-basis simplification stays exact for any $\theta_{14},\theta_{24},\theta_{34}$.
Supplying a neutron density activates
$H_{\rm mat,4}^{\rm flavour}=\diag(V_{CC},0,0,-V_{NC})$
(\cref{ch:core-common}) instead: since this no longer commutes with $O_4$,
the shortcut above is unavailable and \pyfunc{hamiltonian\_matter\_reduced}
takes the general rotated path,
$H_{\rm mat,4}^{\rm red}=O_4^\dagger\,H_{\rm mat,4}^{\rm flavour}\,O_4$,
even when NSI is off.

\textbf{Simultaneous NSI} is fully supported by the same code path
(\cref{sec:nsi-presets}): \pyfunc{NSIConfig.epsilon\_tensor(n\_flavours=4)}
embeds the $3\times3$ $\epsilon$ block into the top-left corner of a
$4\times4$ zero matrix (the sterile row/column stay zero, since $\nu_s$ is
an $\mathrm{SU}(2)_L\times\mathrm{U}(1)_Y$ singlet and carries no NSI
coupling), and \pyfunc{hamiltonian\_matter\_reduced} rotates the resulting
$H_{\rm mat,4}^{\rm NSI,\,flavour}=V\bigl(\diag(1,0,0,0)+\epsilon_4\bigr)$
into the reduced basis with the \emph{same} $O_4^\dagger(\cdot)O_4$
transformation as the 3-flavour case:
\begin{equation}
  H_{\rm mat,4}^{\rm NSI,\,red} = O_4^\dagger\,H_{\rm mat,4}^{\rm NSI,\,flavour}\,O_4,
\end{equation}
again exact for arbitrary $\epsilon$, since it is computed explicitly
rather than assumed to commute with $O_4$ -- the same approximation
trade-off as the 3-flavour NSI case applies here too (previous sections):
the closed-form real-arithmetic shortcut is lost, and diagonalisation falls
back to \code{torch.linalg.eigh}. Nothing in
\pyfunc{hamiltonian\_reduced}/\pyfunc{hamiltonian\_flavour} branches
explicitly on ``sterile or NSI or both'': the flavour count $n\in\{3,4\}$
and the presence of \code{oscillation.nsi} are read directly off
\pyclass{OscillationParameters}, so all four combinations (SM, SM+NSI,
sterile, sterile+NSI) are handled by the exact same two functions.
\end{theorybox}

\begin{notebox}
\textbf{Limitation: $U=O\,U_{\rm red}$ holds exactly only for the
\emph{formal} $U_{\rm red}$, not for \pyfunc{reduced}.} Both here and in
\cref{ch:core-sm}, the abstract contract of \cref{ch:core-common} --
$U=O\,U_{\rm red}$ for the three methods \pyfunc{pmns\_matrix},
\pyfunc{outer\_block}, \pyfunc{reduced} -- is, strictly, a simplification
of what the code returns: \pyfunc{reduced} always drops a $\Delta$ (or
$\Delta_4$) factor that is provably irrelevant for building $H_{\rm kin}$
(any diagonal unitary conjugating $\diag(k)$ leaves it invariant) but whose
omission means \pyfunc{outer\_block}$(\,)\cdot$\pyfunc{reduced}$(\,)$ is
\emph{not} bit-for-bit equal to \pyfunc{pmns\_matrix}$(\,)$. This is
harmless for every use of \pyfunc{reduced} in this project (always as the
sandwiching matrix of a kinetic term), but it means \cref{ch:core-common}'s
statement of the contract should be read as ``$U$ and $O\,U_{\rm red}$
agree up to a diagonal phase that leaves $\diag(k)$ invariant'', not as a
literal matrix identity between the three named methods.
\end{notebox}

\begin{implbox}
\begin{itemize}
  \item \pyclass{PMNSSterileParams}: parameters for \emph{only} the sterile
    extension ($\theta_{14},\theta_{24},\theta_{34}$ and their phases); the
    SM sector lives in a separate \pyclass{PMNSParams} object passed
    alongside this one to the \pyclass{PMNS\_sterile} constructor.
    $\Delta m^2_{41}$ also does not live here: like
    $\Delta m^2_{21}/\Delta m^2_{3\ell}$, it is stored on
    \pyclass{OscillationParameters} because it never enters a mixing
    rotation.
  \item \pyclass{PMNS\_sterile}: subclass of \pyclass{core.common.pmns.PMNS}
    with \code{n\_flavours=4}; it inherits unchanged every piece that is
    generic in size (\pyfunc{R12}, \pyfunc{R13}, \pyfunc{R23},
    \pyfunc{Delta}, \pyfunc{flavour\_basis}, \pyfunc{reduced\_basis},
    \pyfunc{refresh}, \pyfunc{select\_antinu}, ...) and only adds the extra
    constructors \pyfunc{R14}, \pyfunc{R24}, \pyfunc{R34} and the concrete
    implementation of \pyfunc{pmns\_matrix}, \pyfunc{reduced}, and
    \pyfunc{outer\_block}.
  \item Direct consequence of this design: \textbf{every} piece of
    \texttt{medium.*} code that operates on a generic \pyclass{PMNS} object
    works without modification for 4 flavours, simply by passing a
    \pyclass{PMNS\_sterile} instance instead of \pyclass{PMNS\_SM}; the same
    holds for NSI (previous sections), which can be combined freely with
    the sterile extension since \pyfunc{NSIConfig.epsilon\_tensor} embeds
    the $3\times3$ $\epsilon$ block in the top-left corner of the
    $4\times4$ matter matrix, leaving the sterile row/column at zero.
  \item Named presets (e.g.\ \code{"sterile\_3p1\_bestfit\_giunti2017"}) are
    built via \pyfunc{PropagationConfig.oscillation\_parameters\_from\_preset}
    (\cref{ch:core-common}), which internally creates both parameter
    objects and sets $\Delta m^2_{41}$ on the resulting
    \pyclass{OscillationParameters}.
\end{itemize}
\end{implbox}

\begin{refbox}
3+1 mixing parametrisation convention: \cite{KoppMachadoMaltoniSchwetz2013}.
Global 3+1 fit: \cite{GiuntiMarronePalazzo2017}. IceCube atmospheric
$\nu_\mu$ disappearance search: \cite{IceCube2020SterileNumuDisappearance}.
\end{refbox}

\section{Sterile presets}
\label{sec:sterile-presets}

\modulehead{tpeanuts/config/presets.py}{\code{OSCILLATION\_PRESETS} entries
carrying \code{theta14\_deg}: named, literature-justified 3+1 parameter
sets, dispatched to \pyclass{PMNS\_sterile} automatically.}

\begin{implbox}
A preset is a \emph{sterile} preset purely by having a \code{theta14\_deg}
key; \pyfunc{PropagationConfig.oscillation\_parameters\_from\_preset}
(\cref{ch:core-common}) dispatches on that key alone to build a
\pyclass{PMNS\_sterile} instead of a \pyclass{PMNS\_SM}, with every other
field (the SM angles, $\Delta m^2_{21}/\Delta m^2_{3\ell}$) inherited from
the shared NuFIT 5.2 baseline (\cref{ch:core-common}, Default Configuration
box). As with the NSI presets (\cref{sec:nsi-presets}), these are named,
traceable benchmark points rather than a single global best fit -- current
3+1 data is in tension between different anomalies, so no universally
preferred point exists.

\begin{center}
\small
\begin{tabular}{@{}ll@{}}
\toprule
Preset & Parameters and origin \\
\midrule
\parbox[t]{3.2cm}{\code{sterile\_\allowbreak{}3p1\_\allowbreak{}null\_\allowbreak{}mixing}} &
  \parbox[t]{9.4cm}{$\theta_{14}=\theta_{24}=\theta_{34}=0$,
  $\Delta m^2_{41}=1.0\ \mathrm{eV^2}$. Null-hypothesis reference: with
  every sterile angle at zero, \pyclass{PMNS\_sterile} is numerically
  identical to the plain 3-flavour SM (the exact SM-limit embedding proved
  above), run through the 4-flavour machinery -- used to validate that
  embedding in \texttt{core/BSM/test/test3\_bsm\_sterile.py}.} \\[26pt]
\parbox[t]{3.2cm}{\code{sterile\_\allowbreak{}3p1\_\allowbreak{}bestfit\_\allowbreak{}giunti2017}} &
  \parbox[t]{9.4cm}{$\theta_{14}=8.5^\circ$, $\theta_{24}=7.5^\circ$,
  $\theta_{34}=0$, $\Delta m^2_{41}=1.7\ \mathrm{eV^2}$
  ($\sin^2 2\theta_{14}=0.085$, $\sin^2 2\theta_{24}=0.068$). Global 3+1
  fit combining the reactor and gallium anomalies with solar data
  \parencite{GiuntiMarronePalazzo2017}; sterile CP phases fixed to zero
  (no sensitivity in that data combination).} \\[30pt]
\parbox[t]{3.2cm}{\code{sterile\_\allowbreak{}3p1\_\allowbreak{}benchmark\_\allowbreak{}icecube}} &
  \parbox[t]{9.4cm}{$\theta_{14}=0$, $\theta_{24}\approx9.22^\circ$
  ($\sin^2 2\theta_{24}=0.10$), $\theta_{34}=0$,
  $\Delta m^2_{41}=0.3\ \mathrm{eV^2}$. IceCube atmospheric $\nu_\mu$
  disappearance benchmark \parencite{IceCube2020SterileNumuDisappearance};
  $\theta_{14},\theta_{34}$ set to zero since IceCube atmospheric neutrinos
  are primarily sensitive to $\theta_{24}$.} \\[26pt]
\parbox[t]{3.2cm}{\code{sterile\_\allowbreak{}3p1\_\allowbreak{}two\_\allowbreak{}flavour\_\allowbreak{}limit}} &
  \parbox[t]{9.4cm}{$\theta_{12}=\theta_{13}=\theta_{23}=0$,
  $\theta_{14}=15^\circ$, $\theta_{24}=12^\circ$, $\theta_{34}=0$,
  $\Delta m^2_{41}=1.0\ \mathrm{eV^2}$. A mathematical construction, not a
  data-driven benchmark: switching off every SM active angle makes the
  $\nu_e$ disappearance channel reduce exactly to the textbook 2-flavour
  formula $P_{ee}=1-\sin^2(2\theta_{14})\sin^2(1.267\,\Delta m^2_{41}L/E)$
  for any $\theta_{24}$, since $\theta_{24}$ lives only in $O_4$, which
  leaves the electron index untouched. Used to validate the 2-flavour
  approximation against the full 4-flavour evolutor in isolation, before
  quantifying the residual active-sector suppression present for realistic
  presets such as \code{sterile\_3p1\_bestfit\_giunti2017} (see
  \texttt{sterile2\_kinematics.ipynb} in
  \texttt{notebooks/validation/physics/BSM}).} \\[42pt]
\bottomrule
\end{tabular}
\end{center}
\end{implbox}

\begin{notebox}
\textbf{Approximation exposed by \code{sterile\_3p1\_two\_flavour\_limit}.}
The familiar 2-flavour disappearance formula is only \emph{exact} once the
SM active angles are switched off; for realistic presets (SM angles at
their measured values) the $\nu_\mu$ channel in particular picks up a small
residual correction, because $R_{14}$ and $R_{24}$ both act on the sterile
index and do not commute in general -- $|U_{\mu4}|^2=\sin^2\theta_{24}$
holds exactly only when $\theta_{23}=\theta_{14}=0$. This is a property of
the physics (the standard 2-flavour formula is itself an approximation once
more than one active-sterile angle is non-zero), not a code limitation; the
preset exists specifically to isolate and validate this effect
numerically.
\end{notebox}
